\chapter{Wiring}
\label{ch:wiring}

Power in Immersive Engineering travels along wires strung between connectors. The wires sag, they
have length limits, they lose energy over distance, and at high voltage they will kill you. That is
the whole character of the mod in one system.

\section{The three tiers}

\begin{center}
\small
\begin{tabular}{@{}llrrr@{}}
\toprule
Tier & Wire & Max transfer & Loss & Max length \\
\midrule
LV & Copper    & 2\,048 IF/t  & 5\,\%   & 16 blocks \\
MV & Electrum  & 8\,192 IF/t  & 2.5\,\% & 16 blocks \\
HV & Steel     & 32\,768 IF/t & 2.5\,\% & 32 blocks \\
\bottomrule
\end{tabular}
\end{center}

Loss is proportional to how much of the wire's maximum length the run actually uses, so a short HV
run loses very little and a run at the limit loses the full percentage.

\begin{ietip}
Insulated copper and insulated electrum carry the same power as their bare versions and do not shock
anything that touches them. They cost more. On a build people walk through, pay it.
\end{ietip}

\section{Connectors and relays}

A \textbf{connector} is where a wire meets a machine: it takes power from the block it is attached
to, or gives power to it. A \textbf{relay} only passes wire to wire --- it is a pole, not a socket
--- and relays are how you get a line across a distance without a machine at every corner.

\begin{iewarning}
A connector accepts only its own tier's wire. Attaching the wrong one destroys the connector.
\end{iewarning}

\section{Transformers}

Step between tiers. An \textbf{LV/MV Transformer} joins a low-voltage side to a medium-voltage one;
an \textbf{MV/HV Transformer} does the same higher up. Wire the correct tier to each side --- they
are not symmetrical, and the block's model tells you which side is which.

\section{Wire coils and the wirecutter}

Wire is placed by holding a \textbf{coil} and right-clicking one connector, then the other. The
\textbf{Engineer's Wirecutter} removes a wire and returns the coil.

\section{Posts and arms}

\textbf{Wooden} and \textbf{Steel Posts} are what you actually string wire between across country.
Both accept arms, which give you somewhere to mount connectors and transformers off the pole itself.

\section{Breaker switches and meters}

\begin{description}
\item[Breaker Switch] A manual on/off in a wire run, also switchable by redstone.
\item[Redstone Breaker] The same, driven by redstone alone.
\item[Energy Meter] Reports throughput on the line passing through it.
\item[Current Transformer] Emits a redstone signal proportional to the power flowing.
\end{description}

\section{Feedthroughs}

An insulator that carries a wire through a wall without breaking the run, so a line can enter a
building without a connector on each side.

\section{Wire damage}

An energised HV wire damages anything that touches it, and an overloaded wire burns out and drops
as scrap.

\begin{iefork}
City mode changes both. Wire burnout is disabled outright, and shock damage becomes local: rather
than every powered connector broadcasting to the whole network every tick on the chance somebody is
touching a wire, a node that needs the information asks for it. That single broadcast was the most
expensive thing the mod did on a profiled server. \autoref{ch:citymode}.
\end{iefork}

\begin{iefork}
This fork also offers two alternatives to wire entirely: the virtual grid (\autoref{ch:grid}) for
distance, and conduits (\autoref{ch:conduits}) for indoors.
\end{iefork}
